%=======================================================================
%  SUPPLEMENTARY MATERIAL  (compiles as its own PDF)
%=======================================================================
\documentclass[11pt]{article}
\usepackage[a4paper,margin=1in]{geometry}
\usepackage[T1]{fontenc}
\usepackage{lmodern}
\usepackage{amsmath,amssymb}
\usepackage{graphicx}
\graphicspath{{figures/}}
\usepackage{booktabs}
\usepackage{tabularx}
\usepackage{caption}
\usepackage{xcolor}
\definecolor{bibblue}{RGB}{0,63,114}
\usepackage[colorlinks=true,allcolors=bibblue]{hyperref}

\renewcommand{\thetable}{S\arabic{table}}
\renewcommand{\thefigure}{S\arabic{figure}}
\renewcommand{\thesection}{S\arabic{section}}

\title{\bfseries Supplementary Material\\[2pt]
\large How should genomic foundation models be evaluated for metagenomic
taxonomic classification?}
\date{}

\begin{document}
\maketitle

\section{Dataset and splits}
\begin{table}[h]
\centering
\caption{Dataset composition and read-duplication structure.}
\small
\begin{tabularx}{\linewidth}{@{}Xl@{}}
\toprule
Property & Value \\
\midrule
Reference source & Unified human-gut genome catalogue \\
Genera / species & 120 / 1{,}535 \\
Read simulator / length & ART Illumina / 150\,bp \\
Most abundant genus & \emph{Collinsella} (22.8\%) \\
Source pool (total / unique reads) & 258.67M / 42.64M (6.07$\times$) \\
50M training set (unique reads) & 17.76M \\
250M training set (unique reads) & 41.94M ($\approx$ source ceiling) \\
Primary test pool (natural distribution) & 99{,}742 reads, train-disjoint \\
Kraken\,2 coverage-matched pool & 85{,}819 in-database reads \\
\bottomrule
\end{tabularx}
\end{table}

\section{Models and training configuration}
\begin{table}[h]
\centering
\caption{Model configurations. NT-v2 is fine-tuned with LoRA plus an
attention-pooling classification head; MetaTransformer (MT) variants are
trained from scratch and differ only in tokenizer.}
\small
\begin{tabularx}{\linewidth}{@{}lXl@{}}
\toprule
Model & Backbone / tokenizer & Training \\
\midrule
NT-v2 (genus) & Nucleotide Transformer v2, 498M, non-overlap 6-mer & LoRA + attn-pool head \\
NT-v2 shallow & 1-layer random init, non-overlap 6-mer & from scratch \\
MT 6-mer (s1/s6) & MetaTransformer encoder, 6-mer overlap/non-overlap & from scratch \\
MT 12-mer (s1) & MetaTransformer encoder, 12-mer overlap & from scratch \\
MT 13-mer (s1) & MetaTransformer encoder, 13-mer overlap & from scratch \\
\bottomrule
\end{tabularx}
\end{table}

\noindent\textit{[Insert exact hyperparameters before submission: LoRA rank/alpha/dropout,
target modules, optimizer, learning-rate schedule, batch size, epochs, warm-up,
sequence length, and the literal training command for each run. These are
tracked in the \texttt{config.yaml} of each run directory in the code
repository.]}

\section{Full read-level results}
\begin{table}[h]
\centering
\caption{Genus Top-1 accuracy across data scales and tokenizers
(companion to main-text Tables in Sections~4--5). NT-v2 with RC-TTA; MT
forward-only.}
\small
\begin{tabularx}{\linewidth}{@{}Xcc@{}}
\toprule
Setting & Natural pool & Leftover pool \\
\midrule
NT-v2 6-mer 500K            & 55.29\% & --- \\
NT-v2 6-mer 5M             & 63.05\% & --- \\
NT-v2 6-mer 50M            & 67.07\% & 58.1\% \\
NT-v2 6-mer 250M warm       & 67.29\% & 58.3\% \\
NT-v2 6-mer 250M scratch    & 64.8\%  & 55.5\% \\
NT-v2 17.6M species-balanced & 60.4\% & 50.2\% \\
NT-v2 17.6M genus-balanced   & 37.2\% & 35.9\% \\
MT 13-mer 50M              & 87.5\%  & 83.0\% \\
MT 13-mer 250M             & 98.7\%  & 98.5\% \\
MT 12-mer 50M              & 78.83\% & --- \\
MT 6-mer s1 50M            & 48.87\% & 39.1\% \\
MT 6-mer s6 50M            & 47.00\% & 33.0\% \\
\bottomrule
\end{tabularx}
\end{table}

\section{Train-fit ceiling}
\begin{table}[h]
\centering
\caption{Training versus validation accuracy. The 6-mer model cannot fit the
training set past $\sim$66\%; the 13-mer model fits to 99\%.}
\small
\begin{tabularx}{\linewidth}{@{}Xcc@{}}
\toprule
Model & Train acc.\ & Val acc.\ \\
\midrule
NT-v2 250M scratch & 63.2\% & 64.1\% \\
NT-v2 250M warm    & 65.8\% & 66.6\% \\
MT 13-mer 250M     & 98.9\% & 97.5\% \\
\bottomrule
\end{tabularx}
\end{table}

\section{Is 13-mer just lookup? Non-neural baselines}
\begin{table}[h]
\centering
\caption{$k$-mer classifiers on the same 50M-read reference and test pool, with
a 6-mer control that isolates $k$-mer length as the only variable. Genus Top-1
accuracy (\%). Holding the method fixed, naive Bayes on 13-mer ($74.9$) is
$\sim$3$\times$ that on 6-mer ($25.9$); the neural gain is larger on the weaker
6-mer tokenizer ($+41$ pp) than on 13-mer ($+13$ pp); and the parameter-free
13-mer naive Bayes ($74.9$) exceeds the 498M-parameter NT-v2 on 6-mer ($67$).
s1 = overlapping (stride 1); s6 = non-overlapping (stride 6, the NT-v2 setting).}
\small
\begin{tabularx}{\linewidth}{@{}Xccc@{}}
\toprule
Method & 6-mer (s1) & 6-mer (s6) & 13-mer \\
\midrule
Unique-key lookup (Kraken-style) & 0.00 & 0.00 & 0.72 \\
Dominant-genus majority vote & 30.7 & 30.7 & 35.3 \\
Multinomial naive Bayes & 25.7 & 25.9 & 74.9 \\
Neural network & --- & 67.0\textsuperscript{a} & 87.5\textsuperscript{b} \\
\bottomrule
\end{tabularx}

\vspace{2pt}
{\footnotesize \textsuperscript{a}498M-parameter pre-trained NT-v2 + LoRA.\quad
\textsuperscript{b}MetaTransformer, from scratch.}
\end{table}

\section{Computational cost}
\begin{table}[h]
\centering
\caption{End-to-end inference throughput on 100K reads (batch 1024, single
H200) and training peak memory. The large 13-mer vocabulary makes MT 13-mer
slower and more memory-hungry than NT-v2 despite fewer nominal encoder
parameters.}
\small
\begin{tabularx}{\linewidth}{@{}Xccc@{}}
\toprule
Model & reads/s & ms/read & Peak GPU (infer) \\
\midrule
NT-v2 genus       & 1{,}221 & 0.82 & 6{,}731 MiB \\
MT 13-mer s1 genus & 1{,}066 & 0.94 & 9{,}355 MiB \\
MT 6-mer s6        & 3{,}030 & --- & 857 MiB \\
\bottomrule
\end{tabularx}
\end{table}
Training peak memory: MT 13-mer 16.5\,GB vs NT-v2 LoRA 8.8\,GB. Distributed
training on 64 GPUs (8 nodes $\times$ 8 H200) with DDP: 476 min/epoch on a cold
NFS cache versus 43 min/epoch warm (an 11$\times$ I/O effect worth noting for
reproduction at scale).

\end{document}
